\section{Заключение}
\label{sec:Chapter6} \index{Chapter6}

Целью работы было построить математическую модель рекомендательного контура с
переобучением на собственных логах и установить условия его деградации, отделив
вклад переобучения от прочих причин. Вклад работы находится на стыке петель
обратной связи в рекомендательных системах, перформативного предсказания и
теории управления рекомендательными системами~\cite{krauth2022, perdomo2020,
depasquale2026controlSystems}. Поставленные задачи решены полностью.

\textbf{Задача~1.} Рекомендательная система формализована как замкнутая
стохастическая система управления (главы~\ref{sec:Chapter1}
и~\ref{sec:Chapter2}): определены состояние, управляющее воздействие, измерение,
функция перехода~\eqref{eq:closed_loop_transition} и закон адаптации регулятора;
выведено загрязнение логов~\eqref{eq:contaminated} и показано
(предложение~\ref{prop:self_reinforcement}), что обучение на загрязнённых логах
завышает наблюдаемые метрики относительно истинных~--- рост наблюдаемого качества
не равнозначен росту пользы. Этот вывод согласуется с работами о логированных
откликах, алгоритмическом смешении и каузальной коррекции данных
~\cite{sinha2016, chaney2018, bottou2013}.

\textbf{Задача~2.} Доказана Теорема~\ref{thm:beta} о стягивании разнообразия:
внутрикластерный разброс убывает геометрически с множителем $q=1-\beta\rho(2-\beta)$
к асимптотической нижней грани $\beta\sigma_\xi^2/(2-\beta)$. Получены следствия о
времени достижения эхо-камеры и о зависимости границы стягивания от концентрации
выдачи.

\textbf{Задача~3.} Установлено, что стягивание разнообразия универсально по
режиму: оно определяется дрейфом предпочтений и наступает даже без переобучения
(следствия~\ref{cor:no_beta} и~\ref{cor:universal}). Отдельно доказана
Теорема~\ref{thm:loop}: переобучение сужает ковариацию выдачи по сравнению с
замороженной моделью, опуская нижнюю грань разнообразия и переформировывая
ковариацию предпочтений. Тем самым переобучение оставляет отделимый след,
измеримый в дисперсионной части расхождения $\mathrm{KL}(P_T\Vert P_0)$; получены
условия его активации~--- различимость пользователей, наличие сегментной
структуры и порог концентрации.

\textbf{Задача~4.} Предсказания обеих теорем подтверждены серией вычислительных
экспериментов (глава~\ref{sec:Chapter4}), в том числе на наборе
\texttt{MovieLens-20M}: универсальность стягивания, отсутствие эффекта при
$\beta=0$, монотонность по $\alpha\rho$, локализация следа петли в дисперсии и
условие $K\ge2$. Предложен критерий риска и показано (глава~\ref{sec:Chapter5}),
что явная диверсификация выдачи сохраняет разнообразие лучше базовой стратегии
($+53\%$ при том же качестве), тогда как случайное исследование стягивание
усиливает.

Все полученные теоретические результаты строго доказаны и согласуются с
вычислительными экспериментами; центральный тезис работы~--- петля переобучения
оставляет специфический, отделимый от универсального дрейфа след в геометрии
ковариации предпочтений~--- установлен теоретически и подтверждён численно. В
терминах существующей литературы результат уточняет, какой именно компонент
обучения в замкнутом контуре отвечает за наблюдаемую деградацию: динамика с
состоянием даёт универсальное стягивание, а повторное переобучение меняет форму выдачи
~\cite{brown2022stateful, hardt2025performativePastFuture,
barlacchi2026diversityParadox}.

\textbf{Направления дальнейшей работы.} В построенной модели закон отклика
$\theta$ фиксирован, но его текущее значение $\theta_t^{ij}$ меняется вместе с
состоянием пользователя $u_i^t$. Естественным развитием является учёт автономной
динамики предпочтений, не сводящейся к воздействию потреблённых рекомендаций, а
также перенос теоретических оценок с фиксированной когорты на систему с уходом и
приходом пользователей. Перспективны прямое извлечение меры различимости пользователей
$\lVert J_k\rVert$ и резкости выдачи из весов обученной модели как эксплуатационных
метрик риска, раздельный анализ долей следования $s$ и использования $p$, а также
интервенции, нацеленные на форму ковариации выдачи. Отдельного внимания требует
каузальная оценка силы дрейфа $\beta$: полученная в
разделе~\ref{subsec:exp_real} граница по наблюдательным данным не отделяет
рекомендательно-индуцированный дрейф от автономной эволюции вкусов, и для такого
разделения необходима экзогенная вариация показов~--- рандомизованные показы
или регрессия на разрыве у границы отсечения списка. Практическая значимость
результатов состоит в том, что полученные условия допускают мониторинг по
наблюдаемым статистикам и могут служить основой раннего обнаружения деградации
персонализации в промышленных рекомендательных сервисах. Это направление
естественно продолжает каузальные методы для петель обратной связи и безопасные
постановки управления с ограничением на качество базовой политики
~\cite{krauth2022, chow2014baselineControl, chandrasekaran2026networkAware}.
